% Source: Weinberg GR, physical PDF pp. 73--78
%         (printed pp. 47--52).
% Coverage: Section 2.10, Relativistic Hydrodynamics; equations
%           (2.10.1)--(2.10.33), including (2.10.17a).
% Figures/tables/footnotes: no figures or tables; reference marker 4.
% Status: source-reviewed and compile-clean.
% Uncertainties: none.

\subsection{Relativistic Hydrodynamics}\label{sec:2.10}

A great many macroscopic physical systems, including perhaps the universe
itself, may be approximately regarded as \emph{perfect fluids}. A perfect fluid
is defined as having at each point a velocity \(\mathbf v\), such that an
observer moving with this velocity sees the fluid around him as isotropic. This
will be the case if the mean free path between collisions is small compared with
the scale of lengths used by the observer. (For instance, a sound wave will
propagate in air if its wavelength is large compared with the mean free path,
but at very short wavelengths viscosity becomes important and the air stops
acting like a perfect fluid.) We shall translate the above definition of a
perfect fluid into a statement about the energy-momentum tensor.

First suppose that we are in a frame of reference (distinguished by a tilde) in
which the fluid is at rest at some particular position and time. At this
space-time point, the perfect-fluid hypothesis tells us that the energy-momentum
tensor takes the form characteristic of spherical symmetry:
\begin{equation}
  \widetilde T^{ij}=p\,\delta_{ij},
  \tag{2.10.1}\label{eq:2.10.1}
\end{equation}
\begin{equation}
  \widetilde T^{i0}=\widetilde T^{0i}=0,
  \tag{2.10.2}\label{eq:2.10.2}
\end{equation}
\begin{equation}
  \widetilde T^{00}=\rho.
  \tag{2.10.3}\label{eq:2.10.3}
\end{equation}
The coefficients \(p\) and \(\rho\) are called the \emph{pressure} and the
\emph{proper energy density}, respectively. Now go into a reference frame at
rest in the laboratory, and suppose that the fluid in this frame appears to be
moving (at the given space-time point) with velocity \(\mathbf v\). The
connection between the comoving coordinates \(\widetilde x^\mu\) and the lab
coordinates \(x^\mu\) is then
\[
  x^\mu=\Lambda^\mu{}_\nu(\mathbf v)\widetilde x^\nu ,
\]
with \(\Lambda^\mu{}_\nu(\mathbf v)\) the ``boost'' defined by
Eqs.~\eqref{eq:2.1.17}--\eqref{eq:2.1.21}. But \(T^{\mu\nu}\) is a tensor,
so in the lab frame it is
\[
  T^{\mu\nu}
  =\Lambda^\mu{}_\rho(\mathbf v)\Lambda^\nu{}_\sigma(\mathbf v)
   \widetilde T^{\rho\sigma},
\]
or explicitly,
\begin{equation}
  T^{ij}=p\,\delta_{ij}
  +(\rho+p)\frac{v_iv_j}{1-\mathbf v^2},
  \tag{2.10.4}\label{eq:2.10.4}
\end{equation}
\begin{equation}
  T^{i0}=(\rho+p)\frac{v_i}{1-\mathbf v^2},
  \tag{2.10.5}\label{eq:2.10.5}
\end{equation}
\begin{equation}
  T^{00}=\frac{\rho+p\mathbf v^2}{1-\mathbf v^2}.
  \tag{2.10.6}\label{eq:2.10.6}
\end{equation}
To check that this is a tensor, we note that
Eqs.~\eqref{eq:2.10.4}--\eqref{eq:2.10.6} can be integrated into a single
equation:
\begin{equation}
  T^{\mu\nu}=p\eta^{\mu\nu}+(\rho+p)U^\mu U^\nu ,
  \tag{2.10.7}\label{eq:2.10.7}
\end{equation}
where \(U^\mu\) is the velocity four-vector,
\begin{equation}
  \mathbf U=\frac{d\mathbf x}{d\tau}
  =(1-\mathbf v^2)^{-1/2}\mathbf v,\qquad
  U^0=\frac{dt}{d\tau}=(1-\mathbf v^2)^{-1/2},
  \tag{2.10.8}\label{eq:2.10.8}
\end{equation}
normalized so that
\begin{equation}
  U_\mu U^\mu=-1.
  \tag{2.10.9}\label{eq:2.10.9}
\end{equation}
Indeed, Eq.~\eqref{eq:2.10.7} could have been derived very easily by noting
that the quantity on the right-hand side is a tensor, which equals the tensor
\(T^{\mu\nu}\) in a Lorentz frame moving with the fluid, and hence must equal
\(T^{\mu\nu}\) in all Lorentz frames.

Apart from energy and momentum, a fluid will in general carry one or more
conserved quantities, such as the charge, the number of baryons minus the number
of antibaryons, or, at normal temperatures, the number of atoms. Let us
consider one such conserved quantity, and refer to it for brevity as the
``particle number.'' If \(n\) is the particle number density in a Lorentz frame
that moves with the fluid at a given space-time point, then in this frame the
particle-current four-vector at this point is
\begin{equation}
  \widetilde N^i=0,\qquad \widetilde N^0=n .
  \tag{2.10.10}\label{eq:2.10.10}
\end{equation}
In any other Lorentz frame, in which the fluid at this point moves with velocity
\(\mathbf v\), the particle current is related to Eq.~\eqref{eq:2.10.10} by
the ``boost'' \(\Lambda(\mathbf v)\):
\begin{equation}
  N^i=\Lambda^i{}_\mu(\mathbf v)\widetilde N^\mu
  =(1-\mathbf v^2)^{-1/2}v_i n ,
  \tag{2.10.11}\label{eq:2.10.11}
\end{equation}
\begin{equation}
  N^0=\Lambda^0{}_\mu(\mathbf v)\widetilde N^\mu
  =(1-\mathbf v^2)^{-1/2}n ,
  \tag{2.10.12}\label{eq:2.10.12}
\end{equation}
or, more concisely,
\begin{equation}
  N^\mu=nU^\mu .
  \tag{2.10.13}\label{eq:2.10.13}
\end{equation}

The motion of the fluid will be governed by the equations of conservation of
energy and momentum,
\begin{equation}
  0=\partial_\nu T^{\mu\nu}
  =\partial^\mu p
   +\partial_\nu\bigl[(\rho+p)U^\mu U^\nu\bigr],
  \tag{2.10.14}\label{eq:2.10.14}
\end{equation}
and of the particle number:
\begin{equation}
\begin{aligned}
  0=\partial_\mu N^\mu
  &=\partial_\mu(nU^\mu)\\
  &=\frac{\partial}{\partial t}
    \left[n(1-\mathbf v^2)^{-1/2}\right]
   +\boldsymbol{\nabla}\cdot
    \left[n\mathbf v(1-\mathbf v^2)^{-1/2}\right].
\end{aligned}
  \tag{2.10.15}\label{eq:2.10.15}
\end{equation}
It is convenient to write Eq.~\eqref{eq:2.10.14} as separate three-vector and
scalar equations. The three-vector equation is obtained by setting \(\mu=i\)
in Eq.~\eqref{eq:2.10.14}, writing \(U^i=v_iU^0\), and then using
Eq.~\eqref{eq:2.10.14} with \(\mu=0\); this gives
\begin{equation}
  \frac{\partial\mathbf v}{\partial t}
  +(\mathbf v\cdot\boldsymbol{\nabla})\mathbf v
  =-\frac{1-\mathbf v^2}{\rho+p}
  \left[
    \boldsymbol{\nabla}p
    +\mathbf v\frac{\partial p}{\partial t}
  \right].
  \tag{2.10.16}\label{eq:2.10.16}
\end{equation}
The scalar equation is obtained by multiplying Eq.~\eqref{eq:2.10.14} by
\(U_\mu\), using the relation
\begin{equation}
  0=\partial_\nu(U_\mu U^\mu)
  =2U_\mu\partial_\nu U^\mu ;
  \tag{2.10.17}\label{eq:2.10.17}
\end{equation}
we then have
\[
  0=U_\mu\partial_\nu T^{\mu\nu}
  =U^\nu\partial_\nu p
   -\partial_\nu\bigl[(\rho+p)U^\nu\bigr].
\]
Using Eq.~\eqref{eq:2.10.15}, we can write this as
\begin{equation}
\begin{aligned}
  0
  &=U^\nu\left[
      \partial_\nu\rho
      -n\partial_\nu\left(\frac{\rho+p}{n}\right)
    \right]\\
  &=-nU^\nu\left[
      p\,\partial_\nu\left(\frac1n\right)
      +\partial_\nu\left(\frac{\rho}{n}\right)
    \right].
\end{aligned}
  \tag{2.10.17a}\label{eq:2.10.17a}
\end{equation}
The second law of thermodynamics tells us that the pressure \(p\), the energy
density \(\rho\), and the volume per particle \(1/n\) may be expressed as
functions of the temperature \(T\) and the entropy per particle \(\sigma k\)
in such a way that
\begin{equation}
  kT\,d\sigma=p\,d\left(\frac1n\right)
  +d\left(\frac{\rho}{n}\right).
  \tag{2.10.18}\label{eq:2.10.18}
\end{equation}
(Boltzmann's constant \(k\) is introduced here to make \(\sigma\)
dimensionless.) Our scalar equation \eqref{eq:2.10.17a} can now be written
\begin{equation}
  0=U^\mu\partial_\mu\sigma
  \propto\frac{\partial\sigma}{\partial t}
  +(\mathbf v\cdot\boldsymbol{\nabla})\sigma .
  \tag{2.10.19}\label{eq:2.10.19}
\end{equation}
The specific entropy \(\sigma\) is therefore constant in time at any point
that moves along with the fluid. The fundamental equations of relativistic
hydrodynamics are the ``continuity equation'' \eqref{eq:2.10.15}, the ``Euler
equations'' \eqref{eq:2.10.16}, the ``energy equation''
\eqref{eq:2.10.19}, together with equations of state that give \(p\) and
\(\rho\) in terms of \(n\) and \(\sigma\).

In order to gain some insight into the possible equations of state, we may
consider a fluid composed of structureless point particles that interact only in
spatially localized collisions. As shown in Section~\ref{sec:2.8}, the
energy-momentum tensor is
\begin{equation}
  T^{\mu\nu}
  =\sum_N\frac{p_N^\mu p_N^\nu}{E_N}
  \delta^3(\mathbf x-\mathbf x_N).
  \tag{2.10.20}\label{eq:2.10.20}
\end{equation}
[See Eq.~\eqref{eq:2.8.4}.] In a comoving Lorentz frame \(T^{\mu\nu}\) will
have the isotropic form \eqref{eq:2.10.1}--\eqref{eq:2.10.3}, so the
pressure and energy density will be given in this frame by
\begin{equation}
  p=\frac13\sum_{i=1}^{3}T^{ii}
  =\frac13\sum_N\frac{\mathbf p_N^2}{E_N}
    \delta^3(\mathbf x-\mathbf x_N),
  \tag{2.10.21}\label{eq:2.10.21}
\end{equation}
\begin{equation}
  \rho=T^{00}
  =\sum_NE_N\delta^3(\mathbf x-\mathbf x_N),
  \tag{2.10.22}\label{eq:2.10.22}
\end{equation}
whereas the particle number density is, in analogy with Eq.~\eqref{eq:2.6.2},
\begin{equation}
  n=\sum_N\delta^3(\mathbf x-\mathbf x_N).
  \tag{2.10.23}\label{eq:2.10.23}
\end{equation}
It follows that in general
\begin{equation}
  0\leq p\leq\frac{\rho}{3}.
  \tag{2.10.24}\label{eq:2.10.24}
\end{equation}
For a cool, nonrelativistic gas, we can approximate
\[
  E_N\simeq m+\frac{\mathbf p_N^2}{2m},
\]
so Eq.~\eqref{eq:2.10.22} gives
\begin{equation}
  \rho\simeq nm+\frac32p .
  \tag{2.10.25}\label{eq:2.10.25}
\end{equation}
For a hot, extremely relativistic gas, we have
\[
  E_N\simeq\lvert\mathbf p_N\rvert\gg m ,
\]
so Eq.~\eqref{eq:2.10.22} gives
\begin{equation}
  \rho\simeq3p\gg nm .
  \tag{2.10.26}\label{eq:2.10.26}
\end{equation}
Both Eqs.~\eqref{eq:2.10.25} and \eqref{eq:2.10.26} can be incorporated
into a single equation,
\begin{equation}
  \rho-nm\simeq(\gamma-1)^{-1}p ,
  \tag{2.10.27}\label{eq:2.10.27}
\end{equation}
with
\begin{equation}
  \gamma=
  \begin{cases}
    5/3,&\text{nonrelativistic},\\
    4/3,&\text{extreme relativistic}.
  \end{cases}
  \tag{2.10.28}\label{eq:2.10.28}
\end{equation}
Equation \eqref{eq:2.10.18} then gives
\begin{equation}
  kT\,d\sigma
  =p\,d\left(\frac1n\right)
   +(\gamma-1)^{-1}d\left(\frac pn\right)
  =\frac{n^{\gamma-1}}{\gamma-1}
   d\left(\frac{p}{n^\gamma}\right).
  \tag{2.10.29}\label{eq:2.10.29}
\end{equation}
Thus Eq.~\eqref{eq:2.10.19} takes the form
\begin{equation}
  0=\frac{\partial}{\partial t}\left(\frac{p}{n^\gamma}\right)
  +(\mathbf v\cdot\boldsymbol{\nabla})
   \left(\frac{p}{n^\gamma}\right),
  \tag{2.10.30}\label{eq:2.10.30}
\end{equation}
and Eq.~\eqref{eq:2.10.27} is to be used to express \(\rho\) in terms of
\(n\) and \(p\) in Eq.~\eqref{eq:2.10.16}. The proportionality expressed
in Eq.~\eqref{eq:2.10.27} between internal energy and pressure actually holds,
with various values of \(\gamma\), over a class of fluids much wider than the
simple gas of point particles discussed here. For all such fluids, the energy
equation can be put in the form \eqref{eq:2.10.30}.

As an example, let us calculate the speed of sound in a static homogeneous
relativistic fluid. In the unperturbed state, we have \(n,\rho,p,\sigma\)
constant in space and time, and \(\mathbf v=0\). The sound wave produces small
changes \(n_1,\rho_1,p_1,\mathbf v_1\) in \(n,\rho,p,\mathbf v\), but
according to Eq.~\eqref{eq:2.10.19}, it leaves \(\sigma\) unchanged. To first
order in small quantities, Eqs.~\eqref{eq:2.10.15} and
\eqref{eq:2.10.16} then read
\[
  \frac{\partial n_1}{\partial t}
  +n\boldsymbol{\nabla}\cdot\mathbf v_1=0,
  \qquad
  \frac{\partial\mathbf v_1}{\partial t}
  =-\frac{\boldsymbol{\nabla}p_1}{p+\rho}.
\]
But with \(d\sigma=0\), Eq.~\eqref{eq:2.10.18} gives
\[
  0=-\frac{p+\rho}{n}n_1+\rho_1,
\]
so that
\[
  \frac{\partial\mathbf v_1}{\partial t}
  =-\frac{v_s^2\boldsymbol{\nabla}n_1}{n},
\]
where
\begin{equation}
  v_s^2\equiv\frac{p_1}{\rho_1}
  =\left(\frac{\partial p}{\partial\rho}\right)_{\sigma\ \mathrm{const}}.
  \tag{2.10.31}\label{eq:2.10.31}
\end{equation}
Combining the equations for \(n_1\) and \(\mathbf v_1\), we obtain a wave
equation
\[
  0=\left[
    \frac{\partial^2}{\partial t^2}
    -v_s^2\nabla^2
  \right]n_1,
\]
that shows that sound waves travel with the speed \(v_s\), just as in a
nonrelativistic fluid. The speed of sound is much less than the speed of light
(i.e., unity) for a nonrelativistic fluid, but it increases with temperature, so
it is worth checking whether \(v_s\) might exceed unity for a fluid of highly
relativistic point particles, such as hydrogen above \(10^{13}\,{}^\circ\!K\).
In this case, Eqs.~\eqref{eq:2.10.26} and \eqref{eq:2.10.31} give a sound
speed
\begin{equation}
  v_s=\frac1{\sqrt3},
  \tag{2.10.32}\label{eq:2.10.32}
\end{equation}
which is still safely less than unity. This conclusion would not be affected if
electromagnetic forces were taken into account, because Eqs.~\eqref{eq:2.10.7}
and \eqref{eq:2.8.9} impose on the electromagnetic pressure
\(p_{\mathrm{em}}\) and energy density \(\rho_{\mathrm{em}}\) the relation
\begin{equation}
  0=T_{\mathrm{em}}^\mu{}_\mu
  =3p_{\mathrm{em}}-\rho_{\mathrm{em}},
  \tag{2.10.33}\label{eq:2.10.33}
\end{equation}
so the inclusion of \(p_{\mathrm{em}}\) and \(\rho_{\mathrm{em}}\) would not
invalidate Eq.~\eqref{eq:2.10.26} or \eqref{eq:2.10.32}. It is an open
question whether \(v_s\) remains less than unity when nonelectromagnetic forces
are taken into account.\textsuperscript{\hyperref[ch2-ref-4]{4}}
